\scanpage{64}\mspage{6}
% archive sheet number: 66

\[
D\longrightarrow o(D)\qquad (\text{``ombre de }D\text{''})
\]
qui associe à \unsure{tt} disque conforme autour de $D$ \unsure{son} ombre \deleted{\unsure{dans|de}} $T_{X',\unsure{\cdot|a}}$, n'a rien d'injectif. Les \uline{disques ``concentriques''} à $D$ donnent les $\lambda o(D)$ $(0<\lambda\leq1)$, plus précisément, $D\unsure{-|\to}o(D)$ est compatible aux homothéties de rapport \deleted{$\unsure{\lambda>0|\lambda\geq0}$} $\lambda\in\left]0,1\right]$.

Je m'intéresse en fait au cas où $X'$ est compacte, et munie d'une \deleted{\unsure{famille}} partie finie de ``pts marqués'', soit \deleted{$I$}\inserted{$(a_i)_{i\in I}$}, et où on prend une famille de disques $D_i$ $(i\in I)$ \inserted{\uline{mutuellement disjoints}} autour des $a_i$, donnant des ombres
\[
o(D_i)\subset T_{X',a_i}.
\]

On cherche des conditions restrictives \inserted{plausibles} sur le système des $D_i$, impliquant (qui sait?) qu'il est déterminé par le système des ombres.
